\bibitem{TV} B. Toën and M. Vaquié,
\emph{Moduli of objects in dg-categories},
Ann. Sci. École Norm. Sup. (4) \textbf{40} (2007), 387--444.
\url{https://doi.org/10.1016/j.ansens.2007.05.001}.
Theorem 3.6, Corollaries 3.15, 3.17, 3.19, 3.21。
\url{https://www.math.univ-toulouse.fr/~vaquie/marque/vaquie_AnnEns40.pdf}.

\bibitem{ERLowen} W. Lowen,
\emph{Hochschild cohomology, the characteristic morphism and derived deformations},
Compos. Math. \textbf{144} (2008), 1557--1580.
\url{https://arxiv.org/abs/0707.2602v1}.
Theorem 4.8, Section 4.4。任意の小拡大に対するdg加群の記述は本文で説明する。

\bibitem{BFKER} M. Ballard, D. Favero and L. Katzarkov,
\emph{A category of kernels for equivariant factorizations and its implications for Hodge theory},
Publ. Math. Inst. Hautes Études Sci. \textbf{120} (2014), 1--111.
\url{https://doi.org/10.1007/s10240-013-0059-9}.
Theorems 5.39, 6.5。

\bibitem{HRER} D. Huybrechts and J. V. Rennemo,
\emph{Hochschild cohomology versus the Jacobian ring and the Torelli theorem for cubic fourfolds},
Algebraic Geometry \textbf{6} (2019), 76--99.
\url{https://doi.org/10.14231/AG-2019-005}.
Lemma 2.19, Sections 2.5, 3.3。

\bibitem{JLLZER} A. Jacovskis, X. Lin, Z. Liu and S. Zhang,
\emph{Infinitesimal categorical Torelli theorems for Fano threefolds},
J. Pure Appl. Algebra \textbf{227} (2023), 107418.
\url{https://arxiv.org/abs/2203.08187}.
Theorem 4.4(2)。

\bibitem{HuyER} D. Huybrechts,
\emph{The K3 category of a cubic fourfold},
Compos. Math. \textbf{153} (2017), 586--620.
\url{https://doi.org/10.1112/S0010437X16008137}.
Proposition 3.12, Section 3.5。

\bibitem{StacksER} The Stacks Project Authors, \emph{The Stacks Project}.
\url{https://stacks.math.columbia.edu/tag/02HX};
\url{https://stacks.math.columbia.edu/tag/055U}.
2026年11月13日確認。smoothnessの無限小判定とsmooth射のétale局所section。

\bibitem{ERLinZhang} X. Lin and S. Zhang,
\emph{Categorical Torelli and Fourier--Mukai partners of hypersurfaces},
arXiv:2611.05915v1 (2026).
\url{https://arxiv.org/abs/2611.05915}.
Section 2.1, Theorem 2.2, Proposition 2.4（元原稿の参照箇所）。
本稿の判定法の証明の入力ではなく、関連研究として挙げる。

\bibitem{AMSasha} Sasha,
Answer to \emph{Is the property of being a dg generator open?},
MathOverflow, 5 December 2016.
\url{https://mathoverflow.net/questions/256391/is-the-property-of-being-a-dg-generator-open}.
射影直線上の生成対象がskyscraperへ変形して生成性を失う例。
